\section{Research}

The main purpose of this report is testing a wide variety of methods for software development,
in order to get insights to what works or not.
But how do we know what actually works, or will continue to work, and what only happened to work this time?
Like many other fields of study, we try to find the answer through \textit{research}.
But unlike fields such as physics and chemistry, where experiments can be designed to
isolate single variables in controlled conditions, software development involves factors such as domain,
technology as well as differences in group dynamics and individuals,
making it hard to know what is actually affecting the outcome.
This is not a problem specific to software development - Many fields share the difficulties of
minimizing the sheer amount of unknown variables that can affect the results, hindering the ability to make conclusions.
In this section we will explore the role of research in software development by looking at how we can verify hypotheses,
research limitations, how research holds up in practice, and finally we hope to link this up to our own experiences.

\subsection{Research Methods and Limitations}

As mentioned, software development/engineering is not the only field of study where the objective solution is hard to define.
This is common throughout much older fields of studies such as the social-sciences.
In "\textit{The ABC of Software Engineering Research}"\cite{ABC}, Klaas-Jan Stol and Brian
Fitzgerald try and set up a framework in order to make software development/engineering more unified,
largely based on Social Study methodologies.
It highlights that the different methods of conducting research comes with advantages and drawbacks.
An example is case-studies where it highlights strengths such as realistic scenarios and relatively
little researcher interference, but weaknesses such as limited ability to generalize.
Due to these strength and weaknesses, Stol and Fitzgerald argue that to make proper evaluations
in research we should employ different methods so that each can cover for each other's weaknesses.
In the paper "\textit{Design Science in Information Systems Research}"\cite{DesignScience}, Hevner et al.
further highlight strengths and weaknesses of different methods, pointing to a rigor-relevance trade-off.
Striving for scientific precision can require abstracting away variables that play an important role in realistic settings.
The authors also argue that producing and evaluating practical artifacts such as the development methodologies
and practices we explore in this report is itself a rigorous and valuable form of research, not merely an application of theory.

So, what does a software development research paper such as "\textit{A Case Study on the Impact of Refactoring
on Quality and Productivity in an Agile Team}"\cite{Impact_of_Refactoring} tell us?
Alone, it will likely have some of the weaknesses covered by the two papers we
discussed earlier in this section.
Problems with generalization, practical relevance or rigor, and lack of metrics that help us quantify our observations.

The article "\textit{Is it a case study?}"\cite{Is_it_a_case_study} highlights 5 criteria a
study should fulfill to be labeled as a case study: By examining the "\textit{A Case Study on the Impact of Refactoring
on Quality and Productivity in an Agile Team}"\cite{Impact_of_Refactoring} by Moser et.al., we can assess if it qualifies. 
\begin{enumerate}
    \item Empirical investigation of a case
    \item Real-life context
    \item Multiple methods of data collection
    \item Contemporary phenomenon
    \item Will investigators take part
\end{enumerate}

Although the paper is labeled as a case study, it does not fully satisfy the criteria.
In particular, it relies primarily on a single data collection method (automated metrics) and lacks triangulation.
Furthermore, the study is conducted in a “close-to industrial” environment, which weakens its validity as a real-life case study.
Therefore, it is more accurately classified as a small-scale evaluation rather than a true case study.
From this we can see that we should treat the conclusions presented in the paper with caution.
Additionally the choice of Lines of Code is a particularly poor choice of metric.
A reduction in lines of code can be a sign of improved modularity and removal of redundancy.
Therefore, one would likely expect that refactoring would reduce the lines of code in a codebase,
making it an especially poor choice in the context of refactoring.

While there are more than one problem with this article, not all case-studies should be critiqued
by their inability to provide statistical certainty, as that is often not the goal of these studies,
as described in the ABC framework\cite{ABC}.
Rather (when done well), these studies can contribute to the body of knowledge regarding these practices,
and can therefore help us reach conclusions about what methodologies works and under what conditions.

\subsection{Reflections on Our Own Project}

In our own project, we have experimented with frameworks such as design-thinking, agility and scrum
and employed methods such as pair-programming and test-driven development.
Our experiences with these should not be interpreted as truth, but together with
existing articles on these methods: (Evidence for Lean)\cite{Evidence_for_Lean}, (Design thinking was supposed to fix the world)\cite{fix_the_world} \cite{Effectiveness_of_Pair_Programming} (Pair programming),
and (Dark side of Agile)\cite{DarkSide}, we can get a clearer picture on what different methodologies' strengths and weaknesses are.
An example is pair-programming, which is widely considered a very good practice.
We too think that it provided some benefit, but ultimately we found it to be a bit too slow for our schedule,
perhaps due to a difference in experience between our group and many professional developer teams.
Ultimately, our team's experience only contributes a single data point much like the case studies.
Even articles such as "\textit{A Case Study on the Impact of Refactoring on
Quality and Productivity in an Agile Team}"\cite{Impact_of_Refactoring} can be useful to compare
own experiences with others', and also contribute to the larger body of knowledge.
Trying out these practices ourselves also gave an indication of what works or not \textit{for us},
even if that can sometimes conflict with consensus.
